\documentclass[brazil, 12pt]{article}

\usepackage[portuguese]{babel}
\usepackage[utf8]{inputenc}
\usepackage[T1]{fontenc}
\usepackage[scale=0.8]{geometry}
\usepackage{amsmath}
\usepackage{amssymb}
\usepackage{float}
\usepackage{xcolor}
\usepackage{url}
\usepackage{tikz}
\usetikzlibrary{arrows.meta, calc, positioning}

% Cores usadas na figura (para lembrar a caneta azul/vermelha da anotação)
\definecolor{azul}{RGB}{0,80,200}
\definecolor{vermelho}{RGB}{200,30,30}

% Macro para deixar o lugar das figuras indicado sem quebrar a compilação
% caso a imagem ainda não esteja na pasta do Overleaf.
\newcommand{\figplaceholder}[3]{%
\begin{figure}[H]
    \centering
    \IfFileExists{#1}{%
        \includegraphics[width=0.75\linewidth]{#1}%
    }{%
        \fbox{%
        \begin{minipage}[c][4.5cm][c]{0.75\linewidth}
        \centering
        \textbf{Inserir figura aqui}\par
        \vspace{0.3cm}
        Arquivo sugerido: \texttt{#1}
        \end{minipage}}
    }
    \caption{#2}
    \label{#3}
\end{figure}
}

\begin{document}

%-----------------------------------------------------------------------------------------------
%       CABEÇALHO
%-----------------------------------------------------------------------------------------------
\begin{center}
\textbf{Instituto Tecnológico de Aeronáutica - ITA} \\
\textbf{Sistemas de Controle Contínuos e Discretos - CMC-12} \\
\textbf{Aluno}: Matheus Felipe Ramos Borges
\end{center}

\begin{center}
\textbf{Estudo Comparativo de Estimadores de Estado num Pêndulo Invertido sob
Malha de Controle}
\end{center}
%-----------------------------------------------------------------------------------------------
\vspace*{0.5cm}

\section{Introdução}

Este projeto implementa e compara, \textbf{dentro de uma malha de controle
fechada}, diferentes técnicas de \textbf{estimação de estados} aplicadas a um
\textbf{pêndulo invertido sobre carro} (\emph{cart-pole}) --- um sistema
dinâmico não-linear e naturalmente instável.

A motivação parte de uma lacuna do curso: em CMC-12 abordou-se apenas a
\textbf{filtragem passa-baixas clássica}, sem estimação por espaço de estados.
O trabalho avança sobre esse ponto confrontando quatro níveis de estimadores, do
visto em aula ao estado da arte:

\[
\underbrace{\text{Diferenciação} + \text{passa-baixas}}_{\text{\emph{baseline} do curso}}
\;\rightarrow\;
\underbrace{\text{KF linear}}_{\text{linearizado}}
\;\rightarrow\;
\underbrace{\text{EKF}}_{\text{Jacobiano \emph{online}}}
\;\rightarrow\;
\underbrace{\text{UKF}}_{\text{\emph{sigma points}}}.
\]

O pêndulo invertido é escolhido justamente porque o controle opera em
\textbf{duas fases} que exercitam os filtros em regimes distintos: o
\emph{swing-up} (baseado em energia) leva a haste por \textbf{ângulos grandes},
num regime fortemente não-linear onde o KF linear degrada e EKF/UKF se destacam;
já a \emph{estabilização} (LQR na vertical) opera num regime quase-linear, onde
todos os filtros convergem de forma semelhante. Assim, um único experimento
evidencia \textbf{onde cada filtro compensa}.

Quanto ao escopo, o estado é
$X = [\,x\ \ \dot{x}\ \ \theta\ \ \dot{\theta}\,]^{\mathsf T}$ (posição e
velocidade do carro; ângulo e velocidade angular da haste). Apenas as
\textbf{posições} $x$ e $\theta$ são medidas, com ruído (encoders); as
velocidades $\dot{x}$ e $\dot{\theta}$ são estimadas. O desempenho é avaliado por
RMSE por estado, tempo de convergência, robustez (ruído alto, baixa taxa de
amostragem, erro de sintonia de $Q/R$, má inicialização), consistência
estatística (NEES/NIS) e custo computacional, comparando ainda a malha fechada
com estado estimado \emph{vs.} estado ideal (princípio da separação na prática).

\subsection{Parâmetros adotados}

Todas as constantes físicas e de simulação são centralizadas em
\texttt{params.m} (fonte única de verdade) e resumidas na
Tabela~\ref{tab:params}. A haste é modelada como barra uniforme, de modo que o
momento de inércia em torno do centro de massa é $I = \tfrac{1}{3}mL^{2}$.

\begin{table}[H]
\centering
\begin{tabular}{l l l}
\hline
\textbf{Grandeza} & \textbf{Símbolo} & \textbf{Valor} \\
\hline
Massa do carro          & $M$            & $1{,}0$~kg \\
Massa da haste          & $m$            & $0{,}3$~kg \\
Distância pivô--CM       & $L$            & $0{,}5$~m \\
Inércia da haste (CM)   & $I=\tfrac{1}{3}mL^{2}$ & $0{,}025$~kg$\cdot$m$^{2}$ \\
Atrito viscoso do carro & $b$            & $0{,}1$~N$\cdot$s/m \\
Gravidade               & $g$            & $9{,}81$~m/s$^{2}$ \\
\hline
Passo de amostragem     & $\Delta t$     & $0{,}01$~s ($100$~Hz) \\
Horizonte de simulação  & $T_f$          & $10$~s (padrão) \\
\hline
Ruído do encoder ($x$)  & $\sigma_x$     & $0{,}005$~m \\
Ruído do encoder ($\theta$) & $\sigma_\theta$ & $0{,}005$~rad \\
\hline
\end{tabular}
\caption{Parâmetros físicos, de simulação e de medição adotados no projeto.}
\label{tab:params}
\end{table}

\section{Planta de controle}

%-----------------------------------------------------------------------------------------------
%   DIAGRAMA DE CORPO LIVRE NO REFERENCIAL NAO INERCIAL DO CARRINHO
%-----------------------------------------------------------------------------------------------
\noindent{No referencial não inercial do carrinho:}

\vspace{0.4cm}

\begin{center}
\begin{tikzpicture}[
    scale=1.15,
    >=Latex,
    every node/.style={font=\large},
    force/.style={->, thick, line width=1.0pt},
    body/.style={thick, line width=1.1pt},
    rod/.style={line width=3.5pt},
    axis/.style={->, thick, line width=0.9pt},
    measure/.style={<->, thick, line width=0.8pt}
]

    % --- Parametros geometricos -----------------------------------------
    \def\rodang{108}      % angulo da haste medido do eixo x (>90 => inclina para a esquerda)
    \def\rodlen{3.6}      % comprimento desenhado da haste
    \def\comfrac{0.55}    % fracao ate o centro de massa

    % --- Chao -----------------------------------------------------------
    \draw[thick] (-4.6,0) -- (4.8,0);
    \foreach \xx in {-4.4,-3.9,...,4.6}{
        \draw[thin] (\xx,0) -- ++(-0.22,-0.22);
    }

    % --- Carrinho -------------------------------------------------------
    \draw[body, fill=gray!10] (-1.7,0.55) rectangle (1.7,1.45);
    \draw[body, fill=white] (-0.95,0.30) circle (0.24);
    \draw[body, fill=white] (0.95,0.30) circle (0.24);

    % --- Pivo e haste ---------------------------------------------------
    \coordinate (P) at (-0.55,1.45);
    \coordinate (Top) at ($(P)+({\rodlen*cos(\rodang)},{\rodlen*sin(\rodang)})$);
    \coordinate (G)   at ($(P)+({\comfrac*\rodlen*cos(\rodang)},{\comfrac*\rodlen*sin(\rodang)})$);
    \coordinate (V)   at ($(P)+(0,2.9)$);

    \draw[dashed, thick] (P) -- (V);          % vertical de referencia
    \draw[rod] (P) -- (Top);                  % haste
    \fill (P) circle (3.2pt);                 % pivo
    \fill (G) circle (4pt);                   % centro de massa

    % --- Aceleracao angular (theta ponto ponto) -------------------------
    \draw[->, thick] ($(Top)+(0.35,0.05)$) arc[start angle=-20, end angle=60, radius=0.75];
    \node at ($(Top)+(0.15,0.55)$) {$\ddot{\theta}$};

    % --- Pseudo-forca m xddot no CM (aponta para tras, esquerda) --------
    \draw[force] (G) -- ++(-1.25,0);
    \node[anchor=south east] at ($(G)+(-1.25,0.02)$) {$m\ddot{x}$};

    % --- Peso mg ---------------------------------------------------------
    \draw[force] (G) -- ++(0,-1.0);
    \node[anchor=north east] at ($(G)+(-0.05,-0.55)$) {$mg$};

    % --- Reacao vertical no pivo Fy (seta um pouco a direita da tracejada)
    \draw[force, azul] ($(P)+(0.10,0)$) -- ++(0,1.30);
    \node[anchor=west, azul] at ($(P)+(0.22,1.28)$) {$F_y$};

    % --- Comprimento L (rotulo a direita da haste, parte baixa) ---------
    \node[anchor=west] at ($(P)!0.35!(G)+(0.30,0)$) {$L$};

    % --- Angulo theta (entre vertical e haste, junto ao pivo) -----------
    \draw[->, thick] ($(P)+(0,0.50)$) arc[start angle=90, end angle=\rodang, radius=0.50];
    \node at ($(P)+(-0.16,0.34)$) {$\theta$};

    % --- Reacao horizontal no pivo Fx (rotulo abaixo da seta, no carrinho)
    \draw[force, azul] (P) -- ++(-1.05,0);
    \node[anchor=north, azul] at ($(P)+(-0.62,-0.04)$) {$F_x$};

    % --- Forca de controle u (no carrinho, aponta para a direita) -------
    \draw[force, vermelho] (-2.90,0.95) -- (-1.7,0.95);
    \node[vermelho] at (-2.35,1.24) {$u$};

    % --- Atrito viscoso b xdot (embaixo, aponta para a esquerda) --------
    \draw[force, vermelho] (-1.4,0.15) -- ++(-1.1,0);
    \node[anchor=south, vermelho] at (-1.95,0.20) {$b\dot{x}$};

    % --- Termo inercial M xddot (no carrinho, aponta para a esquerda) ---
    \draw[force] (2.75,1.05) -- (1.75,1.05);
    \node[anchor=west] at (2.30,1.35) {$M\ddot{x}$};

    % --- Reacao no carrinho Fx (aponta para a direita) ------------------
    \draw[force, vermelho] (1.75,0.70) -- (2.90,0.70);
    \node[anchor=west, vermelho] at (2.92,0.70) {$F_x$};

\end{tikzpicture}
\end{center}

\vspace{0.4cm}

%-----------------------------------------------------------------------------------------------
%   EQUACOES
%-----------------------------------------------------------------------------------------------

\begin{align}
    u + F_x - b\dot{x} &= M\ddot{x} \\[4pt]
    mgL\sin\theta + m\ddot{x}\,L\cos\theta &= \left(I + mL^{2}\right)\ddot{\theta} \\[4pt]
    F_x + m\ddot{x} &= m\left(L\ddot{\theta}\cos\theta - L\dot{\theta}^{2}\sin\theta\right)
\end{align}

Substituindo $F_x$ da equação (3) na equação (1):

\[
u + m\left(-\ddot{x} + L\ddot{\theta}\cos\theta - L\dot{\theta}^{2}\sin\theta\right) - b\dot{x} = M\ddot{x}
\]

Reorganizando, chega-se ao sistema:

\[
\left\{
\begin{aligned}
u &= (M+m)\ddot{x} + b\dot{x} - mL\ddot{\theta}\cos\theta + mL\dot{\theta}^{2}\sin\theta \\[6pt]
mL\ddot{x}\cos\theta &= \left(I + mL^{2}\right)\ddot{\theta} - mgL\sin\theta
\end{aligned}
\right.
\]

\vspace{0.3cm}

%-----------------------------------------------------------------------------------------------
%   DESENVOLVIMENTO ATE A FORMA DE ESTADOS
%-----------------------------------------------------------------------------------------------

Escolhemos como vetor de estado

\[
X =
\begin{bmatrix}
x_1 \\ x_2 \\ x_3 \\ x_4
\end{bmatrix}
=
\begin{bmatrix}
x \\ \dot{x} \\ \theta \\ \dot{\theta}
\end{bmatrix},
\qquad\text{de modo que}\qquad
\dot{X} =
\begin{bmatrix}
\dot{x} \\ \ddot{x} \\ \dot{\theta} \\ \ddot{\theta}
\end{bmatrix}
=
\begin{bmatrix}
x_2 \\ \ddot{x} \\ x_4 \\ \ddot{\theta}
\end{bmatrix}.
\]

Como $\dot{x}_1 = x_2$ e $\dot{x}_3 = x_4$ são imediatos, resta obter expressões
explícitas para $\ddot{x}$ e $\ddot{\theta}$.


Reescrevendo o sistema com os termos de aceleração à esquerda:

\[
\begin{aligned}
(M+m)\,\ddot{x} \;-\; mL\cos\theta\,\ddot{\theta}
&= u - b\dot{x} - mL\dot{\theta}^{2}\sin\theta \\[4pt]
-\,mL\cos\theta\,\ddot{x} \;+\; \left(I+mL^{2}\right)\ddot{\theta}
&= mgL\sin\theta
\end{aligned}
\]

Podemos usar o formato matricial

\[
\underbrace{
\begin{bmatrix}
M+m & -mL\cos\theta \\[2pt]
-mL\cos\theta & I+mL^{2}
\end{bmatrix}}_{A}
\begin{bmatrix}
\ddot{x} \\[2pt] \ddot{\theta}
\end{bmatrix}
=
\begin{bmatrix}
u - b\dot{x} - mL\dot{\theta}^{2}\sin\theta \\[2pt]
mgL\sin\theta
\end{bmatrix}.
\]

Para simplificar, definimos

\[
a = M+m, \qquad
h = mL\cos\theta, \qquad
J = I + mL^{2},
\]
\[
r_1 = u - b\dot{x} - mL\dot{\theta}^{2}\sin\theta, \qquad
r_2 = mgL\sin\theta,
\]

de forma que $A=\begin{bmatrix} a & -h \\ -h & J \end{bmatrix}$.


O determinante é

\[
\Delta = \det A = aJ - h^{2}
= (M+m)\left(I+mL^{2}\right) - m^{2}L^{2}\cos^{2}\theta,
\]

e a inversa de uma matriz $2\times2$ simétrica é

\[
A^{-1} = \frac{1}{\Delta}
\begin{bmatrix}
J & h \\ h & a
\end{bmatrix}.
\]

Logo,

\[
\begin{bmatrix}
\ddot{x} \\[2pt] \ddot{\theta}
\end{bmatrix}
= \frac{1}{\Delta}
\begin{bmatrix}
J & h \\ h & a
\end{bmatrix}
\begin{bmatrix}
r_1 \\[2pt] r_2
\end{bmatrix}
\quad\Longrightarrow\quad
\begin{aligned}
\ddot{x} &= \frac{J\,r_1 + h\,r_2}{\Delta}, \\[6pt]
\ddot{\theta} &= \frac{h\,r_1 + a\,r_2}{\Delta}.
\end{aligned}
\]


Substituindo $a$, $h$, $J$, $r_1$ e $r_2$:

\[
\ddot{x} =
\frac{
\left(I+mL^{2}\right)\!\left(u - b\dot{x} - mL\dot{\theta}^{2}\sin\theta\right)
+ mL\cos\theta\,\big(mgL\sin\theta\big)
}{
(M+m)\left(I+mL^{2}\right) - m^{2}L^{2}\cos^{2}\theta
},
\]

\[
\ddot{\theta} =
\frac{
mL\cos\theta\,\!\left(u - b\dot{x} - mL\dot{\theta}^{2}\sin\theta\right)
+ (M+m)\,\big(mgL\sin\theta\big)
}{
(M+m)\left(I+mL^{2}\right) - m^{2}L^{2}\cos^{2}\theta
}.
\]

Para a haste uniforme, $I=\dfrac{mL^{2}}{3}$, de modo que $I+mL^{2}=\dfrac{4}{3}mL^{2}$.


Reunindo tudo, com $X=[\,x\ \ \dot{x}\ \ \theta\ \ \dot{\theta}\,]^{\mathsf T}$:

\[
\dot{X} =
\begin{bmatrix}
\dot{x}_1 \\[4pt] \dot{x}_2 \\[4pt] \dot{x}_3 \\[4pt] \dot{x}_4
\end{bmatrix}
=
\begin{bmatrix}
x_2 \\[6pt]
\dfrac{
\left(I+mL^{2}\right)\!\left(u - b\,x_2 - mL\,x_4^{2}\sin x_3\right)
+ m^{2}gL^{2}\sin x_3\cos x_3
}{
(M+m)\left(I+mL^{2}\right) - m^{2}L^{2}\cos^{2} x_3
} \\[16pt]
x_4 \\[6pt]
\dfrac{
mL\cos x_3\,\!\left(u - b\,x_2 - mL\,x_4^{2}\sin x_3\right)
+ (M+m)\,mgL\sin x_3
}{
(M+m)\left(I+mL^{2}\right) - m^{2}L^{2}\cos^{2} x_3
}
\end{bmatrix}.
\]

\vspace{0.3cm}

%-----------------------------------------------------------------------------------------------
%   LINEARIZACAO EM TORNO DA VERTICAL
%-----------------------------------------------------------------------------------------------
% >>> [SEÇÃO B] label adicionado só para permitir referência cruzada a partir da
% nova Seção "Estimadores de Estado Clássicos" (\ref{sec:linearizacao} logo abaixo);
% o conteúdo desta seção não foi alterado.
\section{Linearização em torno da vertical}
\label{sec:linearizacao}

O objetivo do controlador de estabilização é manter a haste na vertical
($\theta = 0$). Linearizamos, portanto, a dinâmica em torno do equilíbrio

\[
X_{eq} = \begin{bmatrix} 0 & 0 & 0 & 0 \end{bmatrix}^{\mathsf T},
\qquad u_{eq} = 0,
\]

usando as aproximações de pequenos ângulos e desprezando os termos de segunda
ordem (o produto $\dot{\theta}^{2}\sin\theta$ é de ordem superior e se anula):

\[
\cos\theta \approx 1, \qquad \sin\theta \approx \theta, \qquad
\dot{\theta}^{2}\sin\theta \approx 0.
\]

Nessas condições o determinante deixa de depender de $\theta$ (sua dependência é
em $\cos^{2}\theta$, cuja derivada em $\theta=0$ é nula):

\[
\Delta_0 = (M+m)\left(I+mL^{2}\right) - m^{2}L^{2}
= mL^{2}\!\left[\tfrac{4}{3}(M+m) - m\right].
\]

As expressões de $\ddot{x}$ e $\ddot{\theta}$ tornam-se lineares:

\[
\begin{aligned}
\ddot{x} &\approx
-\,\frac{\left(I+mL^{2}\right)b}{\Delta_0}\,\dot{x}
\;+\; \frac{m^{2}L^{2}g}{\Delta_0}\,\theta
\;+\; \frac{I+mL^{2}}{\Delta_0}\,u, \\[8pt]
\ddot{\theta} &\approx
-\,\frac{mLb}{\Delta_0}\,\dot{x}
\;+\; \frac{(M+m)mgL}{\Delta_0}\,\theta
\;+\; \frac{mL}{\Delta_0}\,u.
\end{aligned}
\]

Escrevendo $\dot{X} = A\,X + B\,u$, obtemos a \emph{matriz de estados} $A$ e a
\emph{matriz de entrada} $B$:

\[
A =
\begin{bmatrix}
0 & 1 & 0 & 0 \\[4pt]
0 & -\dfrac{\left(I+mL^{2}\right)b}{\Delta_0} & \dfrac{m^{2}L^{2}g}{\Delta_0} & 0 \\[10pt]
0 & 0 & 0 & 1 \\[4pt]
0 & -\dfrac{mLb}{\Delta_0} & \dfrac{(M+m)mgL}{\Delta_0} & 0
\end{bmatrix},
\qquad
B =
\begin{bmatrix}
0 \\[4pt] \dfrac{I+mL^{2}}{\Delta_0} \\[10pt] 0 \\[4pt] \dfrac{mL}{\Delta_0}
\end{bmatrix}.
\]

O termo $A_{4,3} = (M+m)mgL/\Delta_0 > 0$ confirma que a vertical é um equilíbrio
\textbf{instável}: um desvio $\theta>0$ gera $\ddot{\theta}>0$, afastando ainda
mais a haste. É esse comportamento que a realimentação LQR corrige.

\medskip

\noindent\textbf{Instabilidade pelo critério de Routh-Hurwitz.}
Isolando o modo da haste (equação de $\ddot{\theta}$ linearizada), desprezando o
acoplamento com o carro e o atrito e tomando $u=0$, resulta
$\ddot{\theta} - A_{4,3}\,\theta = 0$, cujo polinômio característico é

\[
P(s) = s^{2} + 0\cdot s - A_{4,3},
\qquad A_{4,3} = \frac{(M+m)mgL}{\Delta_0} = 17{,}795 > 0.
\]

A \emph{condição necessária} de estabilidade de Routh-Hurwitz exige que todos os
coeficientes estejam presentes e tenham o mesmo sinal. Aqui ela já é violada em
dois pontos: o coeficiente de $s^{1}$ está ausente ($a_1=0$) e o termo
independente $a_2 = -A_{4,3}$ tem sinal oposto ao de $a_0=1$. Logo, $P(s)$ não
pode ter todas as raízes no semiplano esquerdo.

Montando o arranjo de Routh (com a substituição $a_1 \to \varepsilon \to 0^{+}$
para o zero na primeira coluna):

\[
\begin{array}{c|cc}
s^{2} & 1 & -A_{4,3} \\[2pt]
s^{1} & \varepsilon & 0 \\[2pt]
s^{0} & -A_{4,3} &
\end{array}
\]

A primeira coluna é $\;1,\ \varepsilon,\ -A_{4,3}\;$, isto é, $(+),(+),(-)$:
ocorre \textbf{uma} troca de sinal, indicando \textbf{uma raiz no semiplano
direito}. De fato, as raízes exatas são $s = \pm\sqrt{A_{4,3}} = \pm 4{,}22$, e a
raiz $+4{,}22$ (modo $e^{+4{,}22\,t}$) é a que faz a haste cair. O termo
$-A_{4,3}$ comporta-se como uma ``constante de mola negativa'': é esse polo
instável que a realimentação LQR reposiciona no semiplano esquerdo.

\medskip

\noindent\textbf{Valores numéricos.} Com
$M=1{,}0$~kg, $m=0{,}3$~kg, $L=0{,}5$~m, $g=9{,}81$~m/s$^2$, $b=0{,}1$~N$\cdot$s/m
e $I=mL^{2}/3=0{,}025$~kg$\cdot$m$^2$, tem-se $I+mL^{2}=0{,}1$ e
$\Delta_0 = 0{,}1075$, resultando em

\[
A =
\begin{bmatrix}
0 & 1 & 0 & 0 \\
0 & -0{,}0930 & 2{,}0533 & 0 \\
0 & 0 & 0 & 1 \\
0 & -0{,}1395 & 17{,}795 & 0
\end{bmatrix},
\qquad
B =
\begin{bmatrix}
0 \\ 0{,}9302 \\ 0 \\ 1{,}3953
\end{bmatrix}.
\]

%-----------------------------------------------------------------------------------------------
%   PROJETO DO LQR
%-----------------------------------------------------------------------------------------------
\section{Projeto do controlador LQR}

A escolha do LQR se justifica pela própria natureza do problema. A planta é
\textbf{instável} em malha aberta (Seção anterior, critério de Routh-Hurwitz):
sem realimentação a haste inevitavelmente cai, de modo que um controlador ativo
é obrigatório. Além disso, o cart-pole é um sistema de \textbf{quatro estados
fortemente acoplados com uma única entrada} ($u$, a força no carro); técnicas
clássicas de malha única (como um PID por variável) tornam-se difíceis de
sintonizar nesse cenário multivariável, pois a mesma força precisa regular
simultaneamente posição do carro e ângulo da haste. A realimentação de estados
$u = -K\,X$ ataca todos os estados de uma vez, e o LQR fornece um critério
\textbf{sistemático e ótimo} para escolher os quatro ganhos de $K$: em vez de
posicionar polos por tentativa, minimiza-se um índice de desempenho que
equilibra explicitamente o \emph{erro de regulação} (via $Q$) contra o
\emph{esforço de controle} (via $R$). Como consequência, desde que o par $(A,B)$
seja controlável, o LQR garante malha fechada estável e apresenta boas margens
de robustez. Por fim, o LQR encaixa-se naturalmente no tema deste trabalho: pelo
princípio da separação, o ganho projetado sobre o estado verdadeiro pode operar
sobre o estado \emph{estimado} ($u = -K\,\hat{X}$), permitindo comparar o
desempenho da malha alimentada por cada estimador (KF, EKF, UKF).

\subsection{Controlabilidade}

Antes de projetar, verifica-se que o par $(A,B)$ é controlável, isto é, que a
matriz de controlabilidade tem posto completo:

\[
\mathcal{C} = \begin{bmatrix} B & AB & A^{2}B & A^{3}B \end{bmatrix},
\qquad \operatorname{posto}(\mathcal{C}) = 4.
\]

Para o cart-pole isso se verifica: uma única força aplicada ao carro é capaz de
governar os quatro estados.

\subsection{Formulação do problema}

O regulador linear quadrático escolhe a lei de realimentação de estados
$u = -K\,X$ que minimiza o funcional de custo

\[
J = \int_{0}^{\infty}
\left( X^{\mathsf T} Q\, X + u^{\mathsf T} R\, u \right) dt,
\qquad Q = Q^{\mathsf T} \succeq 0, \quad R = R^{\mathsf T} \succ 0,
\]

em que $Q$ pondera o erro de estado e $R$ pondera o esforço de controle. O ganho
ótimo é

\[
K = R^{-1} B^{\mathsf T} P,
\]

onde $P = P^{\mathsf T} \succeq 0$ é a solução da equação algébrica de Riccati (ARE)

\[
A^{\mathsf T} P + P A - P B R^{-1} B^{\mathsf T} P + Q = 0.
\]

Para um sistema de quarta ordem a ARE não é resolvida à mão; obtém-se $P$ (e daí
$K$) numericamente (Seção~\ref{sec:matlab}).

\subsection{Escolha de $Q$ e $R$ (regra de Bryson)}

Como ponto de partida sistemático, aplica-se a \emph{regra de Bryson}: cada peso
é o inverso do quadrado do desvio máximo aceitável para aquela variável,

\[
Q_{ii} = \frac{1}{x_{i,\max}^{2}}, \qquad R = \frac{1}{u_{\max}^{2}}.
\]

Adotando os limites físicos $x_{\max}=0{,}5$~m, $\dot{x}_{\max}=2$~m/s,
$\theta_{\max}=10^{\circ}\approx0{,}175$~rad, $\dot{\theta}_{\max}=2$~rad/s e
$u_{\max}=20$~N:

\[
Q = \operatorname{diag}(4;\ 0{,}25;\ 32{,}8;\ 0{,}25),
\qquad R = 0{,}0025.
\]

Esses valores são apenas o ponto de partida; a sintonia final é iterativa
(aumentar $Q_{33}$ deixa o ângulo mais firme; aumentar $R$ torna a resposta mais
suave e econômica em força).

\subsection{Cálculo do ganho no MATLAB}
\label{sec:matlab}

Com $A$, $B$, $Q$ e $R$ definidos, o ganho é obtido pela função \texttt{lqr}
(requer a \emph{Control System Toolbox}). O trecho abaixo reaproveita o Jacobiano
analítico já implementado (\texttt{jacobian\_f.m}) para montar $A$:

\begin{verbatim}
p = params();                        % parametros fisicos
A = jacobian_f([0;0;0;0], 0, p);     % matriz de estados (Jacobiano em theta=0)

d     = p.m*p.L^2/3 + p.m*p.L^2;     % I + mL^2
k     = p.m*p.L;                     % mL
a     = p.M + p.m;                   % M + m
Delta = a*d - k^2;                   % determinante em theta = 0
B     = [0; d/Delta; 0; k/Delta];    % matriz de entrada

assert(rank(ctrb(A,B)) == 4, 'par (A,B) nao controlavel');

Q = diag([1/0.5^2, 1/2^2, 1/deg2rad(10)^2, 1/2^2]);
R = 1/20^2;

K = lqr(A, B, Q, R)                  % ganho otimo -> guardar em p.K_lqr
eig(A - B*K)                         % polos de malha fechada (Re < 0 => estavel)
\end{verbatim}

O ganho resultante (valores de referência) é

\[
K = \begin{bmatrix} -40{,}00 & -42{,}14 & 217{,}20 & 48{,}46 \end{bmatrix},
\]

com polos de malha fechada

\[
\begin{aligned}
s_{1,2} &= -12{,}8126 \pm 2{,}5904\,j, \\
s_{3,4} &= -1{,}4435 \pm 1{,}0585\,j,
\end{aligned}
\]

todos com parte real negativa, o que confirma a estabilização da vertical.

%-----------------------------------------------------------------------------------------------
%   SWING-UP ENERGETICO
%-----------------------------------------------------------------------------------------------
\section{Controle de \emph{swing-up} energético}

O LQR só estabiliza a haste dentro de uma pequena vizinhança da vertical (sua
região de atração). Quando o pêndulo parte pendurado, é preciso primeiro
\textbf{levá-lo até essa vizinhança}. Como a linearização não vale para ângulos
grandes, usa-se uma estratégia não-linear baseada em \textbf{energia}
(\emph{energy shaping}), seguindo a formulação de Tedrake~\cite{mit} e o método
clássico de Åström e Furuta~\cite{astrom}: em vez de seguir uma trajetória,
injeta-se energia no pêndulo balançando o carro no ritmo certo, até que ele
tenha energia suficiente para alcançar a vertical.

\subsection{Energia do pêndulo}

Adotamos a notação de~\cite{mit}. A energia da haste é a soma das parcelas
cinética $T$ e potencial $U$,

\[
E(\theta,\dot{\theta}) = T(\dot{\theta}) + U(\theta),
\qquad
T = \tfrac{1}{2} I\,\dot{\theta}^{2},
\qquad
U = -\,m g l\cos\theta,
\]

onde $I$ é o momento de inércia efetivo em torno do pivô e $l$ a distância do
pivô ao centro de massa. Para a haste uniforme deste trabalho,
$I = \tfrac{1}{3}mL^{2} + mL^{2} = \tfrac{4}{3}mL^{2}$ e $l = L$.

\noindent\emph{Observação de convenção.} A referência~\cite{mit} mede $\theta$ a
partir da posição \textbf{pendurada} ($\theta=0$ embaixo), de modo que
$U=-mgl\cos\theta$ é mínimo embaixo. Na modelagem deste relatório $\theta=0$ é a
\textbf{vertical superior}; as duas convenções diferem apenas por uma defasagem
de $\pi$ e um sinal no termo potencial, sem alterar a lei de controle.

A energia da órbita que alcança o equilíbrio superior instável (energia
desejada) é, como em~\cite{mit},

\[
E^{d} = m g l.
\]

O objetivo do \emph{swing-up} é levar $E$ até $E^{d}$. Define-se o erro de energia

\[
\tilde{E} = E - E^{d}.
\]

\subsection{Lei de \emph{energy shaping}}

Para o pêndulo simples atuado por torque na junta, \cite{mit} mostra que a lei

\[
u = -\,k\,\dot{\theta}\,\tilde{E}, \qquad k>0,
\]

produz a dinâmica de erro $\dot{\tilde{E}} = -\,k\,\dot{\theta}^{2}\,\tilde{E}$,
que faz $\tilde{E}\to 0$. Isso decorre da função de Lyapunov
$V = \tfrac{1}{2}\tilde{E}^{2}\ge 0$, cuja derivada é

\[
\dot{V} = \tilde{E}\,\dot{\tilde{E}}
        = -\,k\,\dot{\theta}^{2}\,\tilde{E}^{2} \le 0.
\]

No cart-pole a entrada não é um torque na junta, mas a \textbf{força} $u$ no
carro, que atua sobre a haste através da aceleração do pivô $\ddot{x}$. Da
equação de movimento da haste,
$I\,\ddot{\theta} = m g l\sin\theta + m l\cos\theta\,\ddot{x}$, obtém-se

\[
\boxed{\;\dot{E} = m l\cos\theta\,\dot{\theta}\,\ddot{x}\;}
\]

— os termos gravitacionais se cancelam e a energia só varia por meio de
$\ddot{x}$. Repetindo o argumento de Lyapunov com
$\dot{V} = \tilde{E}\,\dot{E} = \tilde{E}\,m l\cos\theta\,\dot{\theta}\,\ddot{x}$
e escolhendo $\ddot{x}\propto -\,\dot{\theta}\,\tilde{E}\cos\theta$, resulta
$\dot{V} = -\,k\,m l\,\tilde{E}^{2}\,\dot{\theta}^{2}\cos^{2}\theta \le 0$. A lei
de \emph{energy shaping} adaptada ao cart-pole é, portanto, a de~\cite{mit} com
o fator geométrico $\cos\theta$:

\[
\boxed{\;u = -\,k\,\dot{\theta}\,\tilde{E}\,\cos\theta,
\qquad u \leftarrow \operatorname{sat}_{\pm u_{max}}(u).\;}
\]

Quando falta energia ($\tilde{E}<0$) a lei injeta energia
($\dot{E}>0$); ao atingir $E^{d}$ o termo se anula. Nos pontos de retorno
($\dot{\theta}=0$) e na horizontal ($\cos\theta=0$) tem-se $u=0$.

\subsection{Regulação da posição do carro}

A lei de \emph{energy shaping} pura tem uma limitação conhecida: ela governa a
\textbf{energia da haste}, mas não faz nenhuma referência à \textbf{posição do
carro}. Como o bombeamento empurra o carro sempre no sentido que adiciona
energia, o carro deriva sem limite (nas simulações, dezenas de metros), e o LQR
recebe a haste no topo com o carro deslocado --- condição fora de sua região de
atração, de modo que a captura falha.

A solução clássica, seguindo Chung e Spong~\cite{spong} e Tedrake~\cite{mit2},
é somar à lei de \emph{energy shaping} um controlador \textbf{PD sobre o carro},
que o mantém próximo da origem sem impedir o balanço:

\[
\boxed{\;u = -\,k_{swing}\,\dot{\theta}\,\tilde{E}\cos\theta
\;\underbrace{-\,k_p\,x\;-\;k_d\,\dot{x}}_{\text{PD no carro}},
\qquad u \leftarrow \operatorname{sat}_{\pm u_{max}}(u).\;}
\]

O termo $-k_p\,x$ age como uma mola que puxa o carro de volta a $x=0$ e
$-k_d\,\dot{x}$ como amortecimento. Optou-se por um PD (e não PI/PID): o
\emph{swing-up} é um transitório curto sob uma ``perturbação'' oscilatória (o
próprio bombeamento troca de sinal a cada balanço), sem erro de regime
permanente a corrigir --- um termo integral apenas acumularia esse sinal
(\emph{windup}) e brigaria com o bombeamento. Os ganhos adotados, obtidos por
busca em simulação, são

\[
k_{swing} = 12, \qquad k_p = 5, \qquad k_d = 2,
\]

com os quais a haste sobe e é capturada em torno de $1{,}2$~s (estado ideal, sem
ruído). Trata-se de um ponto de partida: os ganhos podem ser reotimizados
(por exemplo, minimizando o tempo de subida sob restrição da excursão do carro).

\subsection{Comutação para o LQR}

O \emph{swing-up} entrega a haste dentro da região de atração do LQR. Adota-se
uma condição de comutação por ângulo e velocidade pequenos,

\[
\lvert\theta\rvert < 15^{\circ}
\quad\text{e}\quad
\lvert\dot{\theta}\rvert < 2~\text{rad/s}
\;\Longrightarrow\;
\text{muda para } u = -K\hat{X},
\]

combinando a fase não-linear (energia) com a fase linear (LQR) num único
controlador de duas fases, exatamente o cenário que exercita os estimadores em
regimes distintos.

%-----------------------------------------------------------------------------------------------
%   MODELO DE SENSORES
%-----------------------------------------------------------------------------------------------
% >>> [SEÇÃO B] label adicionado só para permitir referência cruzada a partir da
% nova Seção "Estimadores de Estado Clássicos"; conteúdo desta seção não foi alterado.
\section{Modelo de sensores e cenários de teste}
\label{sec:sensores}

O modelo de sensores é a interface entre a ``verdade'' do simulador e o que os
estimadores efetivamente observam: recebe o estado verdadeiro e devolve uma
medição imperfeita. Como os filtros nunca acessam $X$ diretamente, é este
modelo que torna o problema de estimação realista.

\subsection{Grandezas medidas}

Apenas as \textbf{posições} são medidas, reproduzindo um par de encoders na
posição do carro e no ângulo da haste; as velocidades $\dot{x}$ e $\dot{\theta}$
não são medidas e devem ser estimadas. A medição é, portanto,

\[
z_k =
\begin{bmatrix} x \\ \theta \end{bmatrix}_{t_k}
+\; v_k,
\qquad
v_k \sim \mathcal{N}(0,\,R),
\qquad
R = \operatorname{diag}\!\left(\sigma_x^{2},\ \sigma_\theta^{2}\right),
\]

com ruído branco gaussiano e independente por canal. Os desvios nominais são
$\sigma_x = \sigma_\theta = 0{,}005$ (Tabela~\ref{tab:params}). Essa escolha ---
medir só posições --- é o que dá sentido à comparação de estimadores: obter
$\dot{x}$ e $\dot{\theta}$ é exatamente a tarefa que separa a simples
diferenciação com passa-baixas dos filtros baseados em modelo (KF/EKF/UKF).

\subsection{Degradações e escolhas de projeto}

Além do ruído nominal, o sensor reproduz três degradações, acionadas por
parâmetros dos cenários (Seção de resultados). As escolhas de modelagem foram:

\begin{itemize}
  \item \textbf{Ruído elevado (S3).} Multiplica-se $\sigma_x$ e $\sigma_\theta$
  por um fator (10$\times$). Testa a filtragem quando a relação sinal-ruído
  cai, favorecendo estimadores que ponderam corretamente medição \emph{vs.}
  modelo.

  \item \textbf{\emph{Dropout} de medida (S4).} A cada passo, com probabilidade
  $p_{drop}$, a leitura é \textbf{totalmente perdida}. Optou-se por sinalizar a
  perda com uma medição \emph{vazia} ($z_k=\varnothing$), em vez de um valor
  qualquer, para que o estimador execute um passo de \textbf{predição pura} (sem
  atualização): a média é propagada pelo modelo e a covariância $P$ cresce,
  refletindo corretamente o aumento de incerteza. Zerar ou congelar a medição
  daria ao filtro uma confiança falsa, corrompendo a análise de consistência
  (NEES/NIS).

  \item \textbf{\emph{Outliers}.} Cada canal, com probabilidade $p_{out}$,
  recebe um pico espúrio de amplitude proporcional a $\sigma$ (por padrão
  $30\sigma$) e sinal aleatório. Modela leituras absurdas e ocasionais, que um
  estimador robusto deve rejeitar. É desligado por padrão
  ($p_{out}=0$) e destina-se a um teste opcional de robustez.
\end{itemize}

Todas as degradações têm probabilidade nula por padrão, de forma que os cenários
sem \emph{override} usam apenas o ruído gaussiano nominal e permanecem
comparáveis entre si.

%=================================================================================================
%   >>> [SEÇÃO B — INÍCIO] ESTIMADORES DE ESTADO CLÁSSICOS
%   Conteúdo redigido a pedido do responsável pela Seção B (README: "KF linear,
%   linearização e filtragem clássica" — arquivos est_lowpass.m, est_complementary.m,
%   est_kf.m). Texto gerado como PONTO DE PARTIDA: revise, reescreva com suas
%   palavras e ajuste os comentários "% >>> EDITAR:" espalhados abaixo antes da
%   entrega final. As equações e os valores numéricos foram conferidos contra o
%   código (src/est_lowpass.m, src/est_complementary.m, src/est_kf.m, src/params.m)
%   e contra um cálculo numérico independente (discretização ZOH), mas vale a pena
%   reconferir depois de qualquer alteração nos parâmetros de params.m.
%=================================================================================================
\section{Estimadores de Estado Clássicos: da Diferenciação ao Filtro de Kalman Linear}
\label{sec:estimadores-classicos}

Esta seção descreve os três estimadores "clássicos" do projeto — isto é, os que
não dependem de linearização \emph{online} nem de amostragem estocástica do
espaço de estados, ao contrário do EKF e do UKF, tratados na seção
seguinte\footnote{% >>> EDITAR: trocar por "Seção~\ref{sec:ekf-ukf}" (ou o label
que for usado) assim que a seção de EKF/UKF do(a) colega responsável pela Seção C
for integrada a este documento — ainda não existe no momento em que este texto
foi escrito.}. Os três compartilham a mesma interface de \emph{drop-in} definida
no contrato do projeto,

\[
[\hat X_k,\ P_k] = \texttt{estimator\_update}(\hat X_{k-1},\, P_{k-1},\, z_k,\, u_k,\, p),
\]

e atuam sobre a mesma medição de posições $z_k=[x\ \ \theta]^{\mathsf T}_{t_k}$
definida na Seção~\ref{sec:sensores}, mas diferem radicalmente na quantidade de
\textbf{modelo} que incorporam: da ausência total de modelo dinâmico (baseline),
passando por um modelo aproximado e implícito (complementar), até o modelo
linear completo do pêndulo, explicitamente propagado com sua incerteza
estatística (Kalman linear). Essa progressão prepara o terreno para o EKF e o
UKF (Seção seguinte), que estendem o mesmo princípio de propagação de
covariância ao modelo \emph{não-linear} do cart-pole.

\subsection{Baseline do curso: diferenciação numérica com filtro passa-baixas}
\label{sec:baseline-lowpass}

O primeiro estimador (\texttt{est\_lowpass.m}) reproduz exatamente a técnica
vista em CMC-12: um filtro passa-baixas digital de primeira ordem (média móvel
exponencial) aplicado à medição, seguido de diferenciação numérica para obter as
velocidades~\cite{oppenheim}. Para cada canal medido,

\[
\tilde x_k = \alpha\,\tilde x_{k-1} + (1-\alpha)\,z_{x,k},
\qquad
\tilde\theta_k = \alpha\,\tilde\theta_{k-1} + (1-\alpha)\,z_{\theta,k},
\qquad \alpha\in(0,1),
\]

e as velocidades são obtidas por diferença atrasada (\emph{backward
difference}) do sinal já filtrado:

\[
\hat{\dot x}_k = \frac{\tilde x_k - \tilde x_{k-1}}{\Delta t},
\qquad
\hat{\dot\theta}_k = \frac{\tilde\theta_k - \tilde\theta_{k-1}}{\Delta t}.
\]

Esse filtro recursivo de primeiro grau é o análogo discreto de um passa-baixas
$RC$ contínuo com constante de tempo $\tau = -\Delta t/\ln\alpha$, ou,
equivalentemente, frequência de corte
$f_c = \dfrac{1}{2\pi\tau} = \dfrac{-\ln\alpha}{2\pi\,\Delta t}$
\cite{oppenheim}. % >>> EDITAR: se preferir, desenvolva a passagem de y[k]=alpha
% y[k-1]+(1-alpha)x[k] até essa fc no corpo do texto (fica como exercício de
% redação — a fórmula em si já está conferida).
Com o valor padrão do código, $\alpha_{lp}=0{,}8$, e $\Delta t = 0{,}01$~s
(Tabela~\ref{tab:params}), obtém-se $f_c \approx 3{,}55$~Hz: qualquer conteúdo do
sinal acima dessa frequência — inclusive parte da dinâmica real do
\emph{swing-up}, que é rápida e fortemente não-linear — é atenuado e atrasado
junto com o ruído do encoder.

Esse é precisamente o ponto fraco do baseline, e a razão de ele servir de
referência inferior na comparação: \textbf{não há nenhum modelo da planta}
embutido no filtro. Quanto maior $\alpha$, menor o ruído mas maior o atraso de
fase (\emph{lag}) — o mesmo compromisso ruído \emph{vs.} atraso de qualquer
filtro passa-baixas clássico~\cite{oppenheim} — e não existe forma de reduzir
simultaneamente os dois, porque o filtro não distingue ruído de medição de
dinâmica real do sistema. Note-se ainda que a covariância $P$ não tem
significado probabilístico aqui: o código apenas a repassa adiante
(\texttt{P = P\_prev}) só para respeitar a interface comum, já que não há
modelo de incerteza a propagar.

\subsection{Filtro complementar}
\label{sec:complementar}

O segundo estimador (\texttt{est\_complementary.m}) já incorpora uma noção
mínima de modelo: a cada passo, prediz a posição por integração de Euler à
frente supondo velocidade constante,

\[
x_{pred} = \hat x_{k-1} + \Delta t\,\hat{\dot x}_{k-1},
\qquad
\theta_{pred} = \hat\theta_{k-1} + \Delta t\,\hat{\dot\theta}_{k-1},
\]

e funde essa predição com a medição bruta usando um ganho fixo
$\alpha_{comp}\in(0,1)$,

\[
x_{fused} = (1-\alpha_{comp})\,x_{pred} + \alpha_{comp}\,z_{x,k},
\qquad
\theta_{fused} = (1-\alpha_{comp})\,\theta_{pred} + \alpha_{comp}\,z_{\theta,k},
\]

com as velocidades novamente obtidas por diferença atrasada do sinal fundido.
O nome "complementar" vem da interpretação clássica no domínio da frequência
introduzida por Higgins~\cite{higgins1975}: dados dois sensores (ou, aqui, um
sensor e um modelo) que medem a mesma grandeza com espectros de ruído
complementares, pondera-se cada fonte por um par de filtros $H(s)$ e $1-H(s)$
que somam a unidade —

\[
\hat X(s) = H(s)\,Z(s) + \big(1-H(s)\big)\,X_{pred}(s),
\]

de modo que a medição (tipicamente confiável em baixa frequência, mas ruidosa
em alta) e o modelo/predição (suave, mas sujeito a deriva em regime
permanente) se complementam sem a necessidade de estimar covariâncias
explicitamente. A versão discreta a ganho fixo implementada aqui é a forma mais
simples dessa ideia; note-se que, com o ganho padrão do código,
$\alpha_{comp} = 0{,}02$, o filtro confia fortemente na predição (98\,\%) e
corrige lentamente pela medição (2\,\%) a cada passo — um compromisso adequado
quando o modelo de velocidade constante é razoável entre amostras, mas que
ainda assim acumula deriva se mantido por muitos passos sem medição (o mesmo
problema clássico de INS/AHRS discutido por Higgins~\cite{higgins1975} e por
Grewal \emph{et al.}~\cite{grewal2013}). É justamente esse ganho fixo,
não-ótimo e invariante no tempo, que o Filtro de Kalman generaliza na próxima
subseção: o KF pode ser visto como um filtro complementar cujo ganho $K_k$ é
recalculado a cada passo para ser o \emph{ótimo} dado o par
modelo/medição~\cite{simon2006, brownhwang2012}.
% >>> EDITAR: se quiser, comente aqui a diferença de comportamento observada
% experimentalmente entre alpha_comp=0,02 (proj. atual) e valores maiores, quando
% os resultados dos cenários S1/S3 estiverem disponíveis.

\subsection{Do modelo linearizado ao Filtro de Kalman discreto}
\label{sec:kf-linear}

\subsubsection{Discretização por segurador de ordem zero (ZOH)}
\label{sec:zoh}

O terceiro estimador (\texttt{est\_kf.m}) usa o modelo \emph{linear}
$\dot X = AX+Bu$ já obtido na Seção~\ref{sec:linearizacao} em torno da vertical,
complementado pelo modelo de medição (linear e exato, já que só se observam
posições diretamente),

\[
z = C X, \qquad
C = \begin{bmatrix} 1 & 0 & 0 & 0 \\ 0 & 0 & 1 & 0 \end{bmatrix}.
\]

Como o KF discreto opera sobre o par $(A_d, B_d)$ amostrado a $\Delta t$, e não
sobre $(A,B)$ contínuos, é preciso discretizar. A discretização exata de um
sistema linear sob entrada constante por trecho (\emph{zero-order hold}, ZOH) é

\[
A_d = e^{A\Delta t}, \qquad
B_d = \int_0^{\Delta t} e^{A\tau}\,d\tau\;B,
\]

resultado padrão de controle digital (ver, por exemplo, Franklin, Powell e
Workman~\cite{franklin1997}, ou Åström e Wittenmark~\cite{astromwittenmark1997}).
Quando $A$ é invertível, $B_d = A^{-1}(A_d - I)B$; \textbf{não é o caso aqui}: a
dinâmica do cart-pole é invariante à translação (a posição $x$ não aparece em
nenhuma equação de movimento), de modo que a primeira coluna de $A$
(Seção~\ref{sec:linearizacao}) é inteiramente nula e $A$ é singular. O caminho
numericamente robusto — e o que o MATLAB usa internamente em \texttt{c2d} — é
formar a exponencial da matriz aumentada~\cite{franklin1997, astromwittenmark1997}

\[
\exp\!\left(
\begin{bmatrix} A & B \\ 0 & 0 \end{bmatrix}\Delta t
\right)
=
\begin{bmatrix} A_d & B_d \\ 0 & I \end{bmatrix},
\]

que evita a inversão e vale mesmo com $A$ singular. Como a saída $z=CX$ é uma
relação puramente algébrica (sem dinâmica), a discretização não altera $C$, ou
seja, $C_d = C$. É exatamente esse procedimento que a função auxiliar
\texttt{linearise\_upright()} (dentro de \texttt{est\_kf.m}) deve executar via
\texttt{c2d(ss(A\_c,B\_c,C\_c,0), dt, `zoh')} — no momento em que este texto foi
escrito, essa função ainda contém \texttt{TODO}s com $A_c$ e $B_c$
\emph{placeholders} (zeros); o preenchimento correto é reaproveitar
literalmente as matrizes $A$ e $B$ já derivadas e verificadas na
Seção~\ref{sec:linearizacao}.
% >>> EDITAR / TODO (código, fora do escopo deste .tex): preencher A_c, B_c em
% est_kf.m com as matrizes de \ref{sec:linearizacao} e chamar linearise_upright()
% a partir de params.m (ou de run_scenario.m) antes de usar est_kf.

Para conferência, com $\Delta t = 0{,}01$~s e os valores numéricos de $A$ e $B$
já apresentados na Seção~\ref{sec:linearizacao}, a discretização ZOH fornece
(calculado numericamente a partir da matriz aumentada acima)

\[
A_d \approx
\begin{bmatrix}
1 & 0{,}009995 & 1{,}03\times10^{-4} & 3{,}4\times10^{-7} \\
0 & 0{,}999070 & 0{,}020529          & 1{,}03\times10^{-4} \\
0 & -6{,}98\times10^{-6} & 1{,}000890 & 0{,}010003 \\
0 & -0{,}001395 & 0{,}177987          & 1{,}000890
\end{bmatrix},
\qquad
B_d \approx
\begin{bmatrix}
4{,}65\times10^{-5} \\ 0{,}009298 \\ 6{,}98\times10^{-5} \\ 0{,}013951
\end{bmatrix}.
\]

Como esperado de uma discretização fiel a um passo pequeno frente à dinâmica
($\Delta t=0{,}01$~s $\ll$ constantes de tempo de $A$, cujos polos ficam em
$\pm4{,}22$ e $-0{,}077$~rad/s — Seção~\ref{sec:linearizacao}), $A_d$ é próxima
da identidade e os autovalores de $A_d$ (isto é, $e^{\lambda_i(A)\Delta t}$)
ficam próximos de $1$: um deles, $e^{+4{,}2266\times0{,}01}\approx1{,}0430$,
maior que $1$ em módulo, confirma que o modo instável da haste
(Seção~\ref{sec:linearizacao}) permanece instável também no domínio discreto,
como deveria.
% >>> EDITAR: se quiser, substitua os valores acima pelos que sua própria
% implementação de linearise_upright()/c2d produzir — devem bater (ZOH é exato).

\subsubsection{Algoritmo do Filtro de Kalman}
\label{sec:kf-algoritmo}

Com $(A_d,B_d,C_d)$ em mãos, o KF linear implementado segue exatamente a forma
padrão de predição--correção de Kalman~\cite{kalman1960, welch1995, simon2006}:

\medskip
\noindent\textbf{Predição} (propaga estado e incerteza pelo modelo, sem usar a medição):
\[
\hat X_{k}^{-} = A_d\,\hat X_{k-1} + B_d\,u_{k-1},
\qquad
P_{k}^{-} = A_d\,P_{k-1}\,A_d^{\mathsf T} + Q,
\]

\noindent\textbf{Atualização} (corrige com a nova medição $z_k$):
\[
S_k = C_d\,P_k^{-}\,C_d^{\mathsf T} + R,
\qquad
K_k = P_k^{-}\,C_d^{\mathsf T}\,S_k^{-1},
\qquad
\hat X_k = \hat X_k^{-} + K_k\big(z_k - C_d\hat X_k^{-}\big),
\qquad
P_k = (I - K_k C_d)\,P_k^{-},
\]

em que $Q$ e $R$ são as covariâncias (supostas conhecidas, diagonais e
constantes) do ruído de processo e de medição, respectivamente
(Tabela~\ref{tab:params-qr} adiante). O ganho $K_k$ não é arbitrário: é o único
ganho que minimiza a incerteza remanescente $\operatorname{tr}(P_k)$ a cada
passo. Partindo da forma de Joseph
$P_k=(I-K_kC_d)P_k^{-}(I-K_kC_d)^{\mathsf T}+K_kRK_k^{\mathsf T}$ (numericamente
mais estável, adotada também no EKF do projeto), derivando
$\operatorname{tr}(P_k)$ em relação a $K_k$ e igualando a zero obtém-se

\[
\frac{\partial}{\partial K_k}\operatorname{tr}(P_k) = 0
\;\Longrightarrow\;
K_k = P_k^{-}C_d^{\mathsf T}\big(C_dP_k^{-}C_d^{\mathsf T}+R\big)^{-1} = P_k^{-}C_d^{\mathsf T}S_k^{-1},
\]

que é exatamente o ganho acima~\cite{kalman1960, simon2006}. Sob as hipóteses
usuais — modelo linear, ruídos brancos, gaussianos, de média nula e
descorrelacionados entre si e do estado inicial — esse ganho torna o KF o
estimador de \textbf{variância mínima entre todos os estimadores lineares não
enviesados} (propriedade BLUE, \emph{Best Linear Unbiased Estimator}), e
coincide com o estimador de máxima verossimilhança/MMSE quando os ruídos são
gaussianos~\cite{kalman1960, simon2006, grewal2013}. É essa otimalidade
comprovada — e não apenas uma escolha heurística de ganho, como no filtro
complementar (Seção~\ref{sec:complementar}) — que justifica o KF como o
patamar de comparação "ótimo" \emph{dentro do regime linear} para o restante do
trabalho.

\paragraph{Escolha de $Q$ e $R$.} O código (\texttt{params.m}) adota, como
ponto de partida,

\begin{table}[H]
\centering
\begin{tabular}{l l}
\hline
\textbf{Matriz} & \textbf{Valor} \\
\hline
$R$ (ruído de medição)  & $\operatorname{diag}(\sigma_x^2,\ \sigma_\theta^2) = \operatorname{diag}(2{,}5\times10^{-5},\ 2{,}5\times10^{-5})$ \\
$Q$ (ruído de processo) & $\operatorname{diag}(10^{-4},\ 10^{-3},\ 10^{-4},\ 10^{-3})$ \\
\hline
\end{tabular}
\caption{Covariâncias de ruído usadas pelo KF linear (compartilhadas com EKF/UKF
via \texttt{p.Q}/\texttt{p.R}, salvo \emph{override} \texttt{p.Q\_kf}/\texttt{p.R\_kf}).}
\label{tab:params-qr}
\end{table}

$R$ decorre diretamente do ruído físico dos encoders (Tabela~\ref{tab:params}) e
não é, a rigor, um parâmetro de sintonia. Já $Q$ modela a incerteza sobre o
quanto o modelo \emph{linearizado} deixa de representar a dinâmica real
(erro de linearização, atrito não modelado, discretização, etc.) e é, na
prática, um parâmetro de ajuste: $Q$ pequeno demais faz o filtro confiar
excessivamente no modelo linear e divergir fora da vizinhança da vertical
(justamente o regime do \emph{swing-up}); $Q$ grande demais aproxima o KF de
uma simples fusão ponderada, perdendo a vantagem de usar o modelo. Vale notar
que o valor de $Q$ aqui foi ajustado diretamente no tempo discreto, e não
obtido pela discretização rigorosa de uma densidade espectral de ruído de
processo contínuo (o que exigiria o método de Van Loan~\cite{vanloan1978}); essa
é uma simplificação comum em projetos aplicados~\cite{simon2006}, mas vale
como ressalva ao interpretar $Q$ como "física" em vez de "efetiva".
% >>> EDITAR: relatar aqui a sensibilidade observada a Q/R nos cenários S5
% (mistuned Q/R) quando os resultados estiverem disponíveis — este parágrafo é
% só a base teórica.

A consistência estatística desse ajuste de $Q$ é o que o teste NEES
(\emph{Normalised Estimation Error Squared}), já implementado em
\texttt{compute\_metrics.m}, mede: sob as hipóteses acima, $e_k^{\mathsf T}P_k^{-1}e_k$
segue uma distribuição $\chi^2$ com 4 graus de liberdade (dimensão do estado); um
NEES sistematicamente acima do intervalo esperado indica $P$
\textbf{otimista} (subestima o erro real — tipicamente $Q$ ou $R$ pequenos
demais), e um NEES sistematicamente abaixo indica $P$ \textbf{pessimista}
(superestima o erro — $Q$ ou $R$ grandes demais)~\cite{barshalom2001}. O mesmo
raciocínio, aplicado à inovação $\nu_k=z_k-C_d\hat X_k^{-}$ em vez do erro de
estado, dá o teste NIS (\emph{Normalised Innovation Squared}, ainda pendente de
implementação em \texttt{compute\_metrics.m} no momento deste texto).

\subsection{Síntese comparativa}
\label{sec:sintese-classicos}

A Tabela~\ref{tab:classicos} resume as diferenças estruturais entre os três
estimadores desta seção.

% >>> [SEÇÃO B] Tabela 3 reformatada: estava com 4 colunas "l" (largura livre)
% e estourava a margem da página porque o cabeçalho da última coluna era uma
% frase longa. Troquei para colunas de largura fixa (p{...}, quebram linha) e
% fonte \small, e encurtei o cabeçalho da última coluna.
\begin{table}[H]
\centering
\small
\begin{tabular}{@{}p{3.1cm} p{3.4cm} p{2.6cm} p{3.3cm}@{}}
\hline
\textbf{Estimador} & \textbf{Modelo dinâmico} & \textbf{Ganho} & \textbf{$P$ tem significado estatístico?} \\
\hline
Baseline (passa-baixas)  & Nenhum                        & Fixo ($\alpha_{lp}$)             & Não \\[2pt]
Complementar             & Implícito (veloc. constante)  & Fixo ($\alpha_{comp}$)           & Não \\[2pt]
KF linear                & $A_d,B_d$ (linearizado)       & Ótimo, variável ($K_k$)          & Sim (sob hipóteses lineares/gaussianas) \\
\hline
\end{tabular}
\caption{Comparação estrutural dos três estimadores clássicos implementados.}
\label{tab:classicos}
\end{table}

% >>> [SEÇÃO B] Esta subseção foi reescrita com resultados reais (antes eram
% apenas expectativas qualitativas + placeholders de figura). Ver
% tests/test_classical_filters.m para o código que gerou os números e as
% figuras abaixo. EDITAR: os parágrafos a seguir descrevem o que os dados
% mostraram — ajuste o texto se reformular o script ou obtiver resultados
% diferentes, e principalmente quando a Fase 2 (harness completo, ver abaixo)
% estiver pronta.

\paragraph{Metodologia (resultado preliminar, pré-Fase~2).} No momento em que
este texto foi escrito, a planta não-linear (\texttt{plant\_cartpole.m}), o
controlador (\texttt{controller\_lqr.m}, \texttt{swingup.m}) e o
\emph{dropout}/\emph{outliers} do sensor (\texttt{sensor\_model.m}) — todos de
responsabilidade da Pessoa A — ainda eram \emph{placeholders}, de modo que o
harness completo (\texttt{main.m}/\texttt{run\_scenario.m}) ainda não roda de
ponta a ponta. Para validar o código desta seção e gerar resultados enquanto
isso, foi escrito \texttt{tests/test\_classical\_filters.m}, que (i)
confere numericamente a linearização de \texttt{linearise\_upright.m} contra o
Jacobiano analítico independente \texttt{jacobian\_f.m} (Pessoa C) — as duas
linearizações, feitas por pessoas diferentes, batem a menos de $10^{-8}$; (ii)
simula uma trajetória próxima ao cenário S1 (haste a $5^{\circ}$ da vertical)
sob realimentação de estado \emph{apenas para fins de demonstração}, usando o
ganho $K$ já publicado na Seção~\ref{sec:matlab} — não é o
\texttt{controller\_lqr.m} definitivo, que continua em aberto; e (iii) roda os
três estimadores desta seção em paralelo sobre essa trajetória, com o ruído
nominal e com o multiplicador $10\times$ do cenário S3. Quando a Fase~2 do
fluxo do projeto (README) estiver pronta, os mesmos gráficos devem ser
regerados a partir de \texttt{run\_scenario.m} com os cenários reais S1/S3, o
que pode alterar os números (mas não deve alterar as conclusões qualitativas
abaixo, já que a única diferença relevante é o controlador de fato usado).

\paragraph{Resultados.} A Tabela~\ref{tab:rmse-classicos} e as
Figuras~\ref{fig:s1-classicos-theta}--\ref{fig:classicos-rmse} resumem o que
foi observado.

\begin{table}[H]
\centering
\small
\begin{tabular}{@{}l r r r r@{}}
\hline
\textbf{Estimador} & \textbf{RMSE $\theta$, S1 (°)} & \textbf{RMSE $\theta$, S3 (°)} & \textbf{RMSE $x$, S1 (m)} & \textbf{RMSE $x$, S3 (m)} \\
\hline
Baseline (passa-baixas) & 0,250 & 1,018 & 0,0031 & 0,0175 \\
Complementar             & 1,036 & 2,500 & 0,0093 & 0,0388 \\
KF linear                & 0,241 & 1,157 & 0,0043 & 0,0188 \\
\hline
\end{tabular}
\caption{RMSE de $\theta$ e de $x$ para os três estimadores clássicos, no
cenário de demonstração descrito acima, com ruído nominal (S1) e com
ruído $10\times$ (S3). Valores gerados por
\texttt{tests/test\_classical\_filters.m}.}
\label{tab:rmse-classicos}
\end{table}

Dois pontos merecem destaque — um deles contraria a expectativa ingênua de que
"mais modelo = sempre melhor":

\begin{itemize}
  \item \textbf{Baseline e KF praticamente empatados, ambos à frente do
  complementar.} No regime quase-linear e bem amostrado ($100$~Hz) deste
  cenário, o passa-baixas bem ajustado ($\alpha_{lp}=0{,}8$) e o KF linear
  entregam RMSE de $\theta$ quase idênticos ($0{,}25^{\circ}$ \emph{vs.}
  $0{,}24^{\circ}$ com ruído nominal). Isso \textbf{não} contradiz a
  otimalidade do KF (Seção~\ref{sec:kf-algoritmo}): o KF é ótimo
  \emph{dado} o par $(Q,R)$ configurado, e neste cenário de demonstração a
  trajetória verdadeira é gerada por uma EDO determinística, sem ruído de
  processo de fato injetado — ou seja, o $Q$ de \texttt{params.m}
  (Tabela~\ref{tab:params-qr}), calibrado para representar erro de
  linearização e dinâmica não modelada, é relativamente "generoso" frente ao
  ruído de processo real (próximo de zero) deste cenário específico. Isso dá
  ao ganho do KF uma folga maior do que o estritamente necessário, deixando
  passar um pouco mais de ruído de medição do que um passa-baixas
  fortemente suavizante deixaria — o mesmo fenômeno de sensibilidade a
  $Q/R$ mistunado que o cenário S5 é desenhado para expor
  (Seção~\ref{sec:kf-algoritmo}). A vantagem estrutural do KF (explorar
  $A_d,B_d$) deve ficar mais evidente quando houver ruído de processo real
  (p.ex. erro de linearização acumulado fora da vizinhança da vertical) ou
  quando $Q$ for recalibrado — ambos fora do escopo desta seção.

  \item \textbf{Filtro complementar é o pior dos três, por larga margem.} A
  Figura~\ref{fig:s1-classicos-theta} mostra o motivo: com o ganho padrão
  $\alpha_{comp}=0{,}02$ (Seção~\ref{sec:complementar}), o filtro confia
  $98\,\%$ na predição de velocidade constante a cada passo e corrige
  lentamente pela medição — adequado para deriva lenta, mas lento demais
  para acompanhar a oscilação inicial rápida de $5^{\circ}$ deste cenário: o
  filtro \emph{overshoot}a a mais de $6^{\circ}$ e só converge por volta de
  $t\approx2$~s, enquanto baseline e KF acompanham a verdade quase
  imediatamente. Sob ruído $10\times$ (S3), seu RMSE mais que dobra
  ($1{,}04^{\circ}\to2{,}50^{\circ}$), a maior degradação relativa das três
  — esperado, já que $\alpha_{comp}$ é fixo e não se adapta ao novo nível de
  ruído, ao contrário do ganho $K_k$ do KF, recalculado a cada passo.
\end{itemize}

Nenhum dos três estimadores desta seção foi testado no \emph{swing-up}
(cenário~S2, ângulos grandes): isso exigiria o controlador de \emph{swing-up}
(Pessoa A, ainda TODO) e, mais importante, é precisamente o regime em que a
linearização por trás do KF deixa de valer — razão pela qual o EKF/UKF (Seção
seguinte) são necessários, e não uma limitação deste teste específico.

\figplaceholder{results/S1\_classicos\_theta.png}{Ângulo $\theta$ estimado
pelos três estimadores clássicos no cenário de demonstração (análogo a S1,
$5^{\circ}$ de perturbação inicial, ruído nominal) \emph{vs.} estado verdadeiro
simulado. O filtro complementar (ganho fixo $\alpha_{comp}=0{,}02$)
\emph{overshoot}a e leva ${\sim}2$~s para convergir; baseline e KF acompanham a
verdade quase imediatamente.}{fig:s1-classicos-theta}

\figplaceholder{results/S1\_S3\_classicos\_rmse.png}{RMSE de $\theta$ dos três
estimadores clássicos no cenário de demonstração, com ruído nominal e com
ruído $10\times$ (análogo a S1 \emph{vs.} S3). O complementar tem a maior
degradação relativa; baseline e KF permanecem próximos em ambos os
níveis.}{fig:classicos-rmse}

%=================================================================================================
%   >>> [SEÇÃO B — FIM]
%=================================================================================================

%-----------------------------------------------------------------------------------------------
%   DISCUSSAO / ARQUITETURA DA MALHA
%-----------------------------------------------------------------------------------------------
\section{Arquitetura da malha e discussão}

A Figura~\ref{fig:malha} reúne os módulos num diagrama de blocos da malha
fechada. O laço opera em tempo discreto (passo $\Delta t$): a cada instante o
\textbf{sensor} mede as posições ($z=[x;\theta]$ com ruído/perda), o
\textbf{estimador} (KF, EKF ou UKF) reconstrói o estado completo $\hat{X}$, o
\textbf{controlador} de duas fases gera a força $u$ e a \textbf{planta}
não-linear é integrada até o passo seguinte.

\begin{figure}[H]
\centering
\begin{tikzpicture}[
  >=Latex,
  node distance=1.5cm and 1.9cm,
  block/.style={draw, thick, rounded corners, align=center,
                minimum height=1.25cm, minimum width=2.5cm},
  sig/.style={->, thick, line width=0.9pt}
]
  \node[block] (ctrl) {Controlador\\[-2pt]\small swing-up + PD\\[-2pt]\small $\rightarrow$ LQR};
  \node[block, right=of ctrl] (plant) {Planta\\[-2pt]\small cart-pole\\[-2pt]\small (não-linear)};
  \node[block, right=of plant] (sensor) {Sensor\\[-2pt]\small ruído, dropout,\\[-2pt]\small outliers};
  \node[block, below=1.6cm of sensor] (est) {Estimador\\[-2pt]\small KF / EKF / UKF};

  \draw[sig] ($(ctrl.west)+(-1.5,0)$) -- node[above,align=center]{$X^\ast{=}0$}(ctrl.west);
  \draw[sig] (ctrl) -- node[above]{$u$} (plant);
  \draw[sig] (plant) -- node[above,align=center]{$x$\\\small (verdade)} (sensor);
  \draw[sig] (sensor) -- node[right]{$z$} (est);
  \coordinate (fb) at (ctrl.south |- est.west);
  \draw[sig] (est.west) -- node[below]{$\hat{X}$ (estimado)} (fb) -- (ctrl.south);
\end{tikzpicture}
\caption{Malha de controle com estimação de estados. O controlador nunca acessa
o estado verdadeiro: fecha o laço sobre a estimativa $\hat{X}$ (princípio da
separação). É neste ponto que se comparam os quatro estimadores.}
\label{fig:malha}
\end{figure}

Dois aspectos merecem destaque na discussão. Primeiro, o
\textbf{princípio da separação}: o controlador foi projetado supondo acesso ao
estado, mas na prática opera sobre $\hat{X}$; a qualidade de cada estimador
reflete-se diretamente no desempenho da malha, e é exatamente essa degradação
que o trabalho quantifica. Segundo, o \textbf{controlador de duas fases}: longe
da vertical usa a lei não-linear de \emph{energy shaping} com o PD de carro;
dentro da janela de captura comuta para o LQR linear. Assim, um único
experimento leva os filtros por um regime fortemente não-linear (\emph{swing-up})
e por um regime quase-linear (estabilização), evidenciando onde cada técnica de
estimação compensa.

%-----------------------------------------------------------------------------------------------
%   REFERENCIAS
%-----------------------------------------------------------------------------------------------
% >>> [SEÇÃO B] argumento trocado de {9} para {99}: é só a largura reservada
% para o rótulo numérico (afeta a indentação), e passamos de 5 para 13
% referências com a adição das entradas da Seção B logo abaixo.
\begin{thebibliography}{99}

\bibitem{mathworks}
MathWorks. \emph{Derive and Simulate Cart-Pole System} (Symbolic Math Toolbox).
Disponível em:
\url{https://www.mathworks.com/help/symbolic/derive-and-simulate-cart-pole-system.html}.

\bibitem{mit}
R.~Tedrake. \emph{Underactuated Robotics: Algorithms for Walking, Running,
Swimming, Flying, and Manipulation} --- Cap. \emph{The Simple Pendulum}.
Disponível em: \url{https://underactuated.mit.edu/pend.html}.

\bibitem{astrom}
K.~J.~Åström, K.~Furuta. \emph{Swinging up a Pendulum by Energy Control}.
Automatica, vol.~36, n.~2, pp.~287--295, 2000.

\bibitem{mit2}
R.~Tedrake. \emph{Underactuated Robotics} --- Cap. \emph{Acrobots, Cart-Poles,
and Quadrotors}. Disponível em: \url{https://underactuated.mit.edu/acrobot.html}.

\bibitem{spong}
C.~C.~Chung, J.~Hauser. \emph{Nonlinear control of a swinging pendulum}.
Automatica, vol.~31, n.~6, pp.~851--862, 1995. Ver também
M.~W.~Spong, \emph{Energy Based Control of a Class of Underactuated Mechanical
Systems}, IFAC World Congress, 1996.

% >>> [SEÇÃO B — INÍCIO DAS REFERÊNCIAS ADICIONADAS] usadas em
% Seção~\ref{sec:estimadores-classicos} (baseline, filtro complementar, KF linear
% e discretização ZOH). Vale a pena consultar estas na ordem: (1) Kalman 1960 é o
% artigo original (curto, histórico, mais difícil de ler hoje em dia); (2) Welch
% & Bishop é o tutorial mais didático para começar; (3) Simon 2006 e Grewal &
% Andrews são os livros-texto de referência, com os capítulos de KF/EKF/UKF todos
% num só lugar; (4) Higgins 1975 é a referência clássica do filtro complementar;
% (5) Franklin/Powell/Workman e Åström & Wittenmark cobrem a discretização ZOH
% (conteúdo de controle digital, alinhado à ementa de CMC-12); (6) Van Loan 1978
% e Bar-Shalom et al. são mais avançados/opcionais (ruído de processo contínuo
% "correto" e testes de consistência NEES/NIS, respectivamente).
\bibitem{kalman1960}
R.~E.~Kalman. \emph{A New Approach to Linear Filtering and Prediction
Problems}. Journal of Basic Engineering (ASME Transactions), vol.~82, Series
D, pp.~35--45, 1960.

\bibitem{welch1995}
G.~Welch, G.~Bishop. \emph{An Introduction to the Kalman Filter}. Technical
Report TR~95-041, University of North Carolina at Chapel Hill, 1995
(atualizado em 2006). Disponível em:
\url{https://www.cs.unc.edu/~welch/media/pdf/kalman_intro.pdf}.

\bibitem{simon2006}
D.~Simon. \emph{Optimal State Estimation: Kalman, $H_\infty$, and Nonlinear
Approaches}. Wiley, 2006.

\bibitem{grewal2013}
M.~S.~Grewal, A.~P.~Andrews. \emph{Kalman Filtering: Theory and Practice Using
MATLAB}, 4.ª ed. Wiley, 2015 (1.ª ed. 1993).

\bibitem{higgins1975}
W.~T.~Higgins. \emph{A Comparison of Complementary and Kalman Filtering}. IEEE
Transactions on Aerospace and Electronic Systems, vol.~AES-11, n.~3,
pp.~321--325, 1975.

\bibitem{brownhwang2012}
R.~G.~Brown, P.~Y.~C.~Hwang. \emph{Introduction to Random Signals and Applied
Kalman Filtering}, 4.ª ed. Wiley, 2012.

\bibitem{franklin1997}
G.~F.~Franklin, J.~D.~Powell, M.~L.~Workman. \emph{Digital Control of Dynamic
Systems}, 3.ª ed. Addison-Wesley, 1997.

\bibitem{astromwittenmark1997}
K.~J.~Åström, B.~Wittenmark. \emph{Computer-Controlled Systems: Theory and
Design}, 3.ª ed. Prentice Hall, 1997.

\bibitem{oppenheim}
A.~V.~Oppenheim, R.~W.~Schafer. \emph{Discrete-Time Signal Processing}, 3.ª
ed. Prentice Hall, 2009.

\bibitem{vanloan1978}
C.~F.~Van Loan. \emph{Computing Integrals Involving the Matrix Exponential}.
IEEE Transactions on Automatic Control, vol.~23, n.~3, pp.~395--404, 1978.

\bibitem{barshalom2001}
Y.~Bar-Shalom, X.~R.~Li, T.~Kirubarajan. \emph{Estimation with Applications to
Tracking and Navigation: Theory, Algorithms and Software}. Wiley, 2001.
% >>> [SEÇÃO B — FIM DAS REFERÊNCIAS ADICIONADAS]

\end{thebibliography}

\end{document}